\documentclass[11pt,a4paper]{article}
% Estilo comum dos PDFs individuais da Sprint 1.
\usepackage[utf8]{inputenc}
\usepackage[T1]{fontenc}
\usepackage[brazilian]{babel}
\usepackage[a4paper,margin=2.35cm]{geometry}
\usepackage{xcolor}
\usepackage{hyperref}
\usepackage{enumitem}
\usepackage{listings}
\usepackage{parskip}
\usepackage{needspace}
\definecolor{azul}{RGB}{20,63,112}
\hypersetup{colorlinks=true,linkcolor=azul,urlcolor=azul}
\setlist[itemize]{leftmargin=*,itemsep=3pt}
\lstset{basicstyle=\ttfamily\small,frame=single,breaklines=true,columns=fullflexible}
\sloppy
\newcommand{\codigo}[1]{\texttt{#1}}
\newcommand{\titulo}[2]{\begin{center}{\LARGE\bfseries #1\par}\vspace{4pt}{\large #2\par}\vspace{8pt}\rule{0.8\textwidth}{0.4pt}\par\vspace{5pt}{\small Matheus Moreli Durão Silva -- Matrícula 12621CBS217\par Cibersegurança - Bacharelado - Noturno\par}\end{center}}
\newcommand{\cardfonte}[1]{\par\smallskip\noindent\textit{Referência: }\href{#1}{Trello - enunciado e checklist do card}.}

\begin{document}\titulo{S1\_20 - Revisão geral}{Conectar cada etapa ao sistema final}
\section*{Finalidade deste card}
Revisar não é repetir o nome dos assuntos. É conseguir olhar o \codigo{app.py}, explicar qual recurso Python resolve cada problema e conferir se os exercícios e entregáveis foram enviados. Este recorte organiza a revisão da Sprint 1 e registra os pontos a retestar.
\section*{Mapa de evolução}
S1\_01 a S1\_03 definem problema, requisitos e testes. S1\_04 a S1\_09 ensinam saída, dados, decisões, repetição e funções: a base do menu. S1\_10 a S1\_13 estruturam informações com listas, sets, dicionários e enumerações. S1\_14 e S1\_15 fazem os dados sobreviverem ao programa e tratam falhas. S1\_16 a S1\_19 registram essa evolução com Git, branches e PRs.
\section*{Perguntas de autoavaliação}
\begin{itemize}
\item Consigo cadastrar um ativo e apontar onde cada entrada é validada?
\item Consigo explicar por que ativos usam dicionários, vulnerabilidades usam lista associada e JSON salva tipos simples?
\item Consigo diferenciar \codigo{if}, \codigo{while}, \codigo{for}, função, parâmetro e \codigo{return} com um exemplo próprio?
\item Consigo mostrar no Git qual commit/PR levou uma funcionalidade à \codigo{main}?
\end{itemize}
\section*{Como confirmar entregas}
Abrir cada pasta S1\_01 a S1\_20, ler o enunciado e checklist correspondente, executar o arquivo Python quando houver um pedido de código e abrir o PDF como evidência explicada. O status no Trello deve ser conferido separadamente: este documento não altera cards nem declara concluído o que ainda está em andamento.
\section*{Próximo uso}
Se uma pergunta não puder ser respondida com exemplo e teste, ela vira uma dúvida de estudo. A Sprint 2 deve aprimorar o que existe, não esconder conceitos da Sprint 1. Por isso, cada entrega individual mantém o enunciado, a explicação, a aplicação e a forma de conferência.
\cardfonte{https://trello.com/c/X3LzSoi0}
\end{document}
